% Secao 6 - o coracao do relatorio: concordancia e divergencia distancia x fluxo.
% Disciplina: descritivo, associativo. Numeros = macros de values.tex.

\section{Concordância e divergência entre os dois desertos}
\label{sec:divergencia}

As seções anteriores construíram duas medidas de deserto hospitalar: a distância
ao hospital ativo mais próximo, que uma política de acesso baseada em mapa
enxerga, e o burden efetivo de fluxo F1, que pondera a viagem por onde os
pacientes de fato internam. Se as duas dissessem a mesma coisa, F1 seria
redundante e este relatório, desnecessário. Mas não dizem: concordam o
suficiente para medir o mesmo fenômeno e divergem o suficiente para que a escolha
da medida mude quem é classificado como desassistido. Resta caracterizar onde, e
para quem, a divergência aparece.

\subsection{Quanto um conceito é redundante com o outro}

A Figura~\ref{fig:scatter} cruza as duas medidas no município. A correlação é
moderada: $r = \valCorrFD$, de modo que a distância ao hospital mais próximo
explica cerca de \valRtwoFD\% da variação no burden efetivo de fluxo.% src: 11_fig2_scatter; corr 0.551, R2 0.304
O que sobra não é ruído. Quase toda a massa de pontos está \emph{acima} da
diagonal $F1 = \text{distância}$: o paciente típico não interna no hospital mais
próximo, e o burden efetivo mediano (\valMedF~km) é o triplo da distância mediana
ao hospital mais próximo (\valMedIso~km), uma diferença mediana de \valMedGap~km
que mede o desvio revelado em relação ao mapa.% src: 01_build_f1_divergence headline
A consequência para mensuração é direta. Definindo deserto como o quartil
superior de cada medida, dos municípios sinalizados como deserto por
\emph{algum} dos dois conceitos, \valPctDisagree\% o são por apenas um deles; e
\valPctFlowMissed\% dos desertos de fluxo são invisíveis ao conceito de
distância.% src: bloco síntese - pct_disagree 60.5, pct_flow_missed 43.3
Trocar quilômetros por fluxo revelado não recalibra uma medida --- redesenha o
mapa de quem precisa de atenção.

\subsection{O mapa da divergência}

A Figura~\ref{fig:divergencia} pinta cada município pelo quadrante que ocupa, e
a geografia do desacordo não é aleatória. Os desertos de fluxo ocultos à
distância (vermelho) concentram-se no interior do Nordeste e do Centro-Oeste:
municípios com um hospital por perto no mapa, mas cuja necessidade de internação
escoa para hubs distantes. O padrão inverso (azul) domina a Amazônia, onde a
distância rodoviária ao hospital mais próximo é grande mas o fluxo efetivo é
curto --- reflexo de que a malha de estradas não representa o deslocamento real,
por rio, e de que o atendimento, quando ocorre, ocorre localmente. O deserto
pleno, em que os dois conceitos concordam, ocupa o Norte mais remoto; a
concordância de que não há deserto cobre o Sul e o Sudeste populosos.

\subsection{Quem vive em cada quadrante}

A Tabela~\ref{tab:perfis} traça o perfil dos quadrantes e revela três padrões.
O primeiro desfaz o paradoxo aparente do quadrante perto-mas-isolado: longe de
não terem oferta, esses municípios têm a \emph{maior} presença de hospital local
da tabela (\valFOhosp\% têm ao menos um leito hospitalar instalado) e ainda assim
90\% de suas internações saem do município.% src: profiles_by_quadrant.csv
O hospital local existe, mas resolve pouco; a necessidade real --- a que exige
complexidade --- viaja. É exatamente o tipo de isolamento que a contagem de
leitos no mapa não captura e o fluxo revela.

O segundo padrão é distributivo. Os três quadrantes de deserto são
sistematicamente mais pobres (PIB per capita em torno de R\$~13--14~mil) do que o
quadrante conectado (R\$~22~mil), de modo que a fricção de acesso, medida por
qualquer dos conceitos, recai sobre os municípios de menor renda. O terceiro
padrão qualifica o papel do setor privado: a cobertura de saúde suplementar é
baixa nos desertos (\valBOans\% a \valDOans\%) e bem maior no quadrante conectado
(\valNEans\%). Onde o fluxo revela isolamento, não há rede privada que absorva a
demanda --- o gap morde com força em municípios dependentes do SUS. (A relação
dos quadrantes com mortalidade evitável é deliberadamente adiada para a
Seção~\ref{sec:mortalidade}, que mostra por que a associação simples engana.)

\begin{table}[t]
\centering
\caption{Perfil dos quadrantes distância $\times$ fluxo}
\label{tab:perfis}
\footnotesize
\begin{tabular}{lrrrrrrr}
\toprule
Quadrante & Munic. & Pop. & Dist. & F1 & Fluxo que & Hosp. & PIB pc \\
          &        & (mi) & (km)  & (km) & sai      & local & (R\$ mil) \\
\midrule
Perto, deserto em fluxo & \valFOn & \valFOpop  & \valFOiso & \valFOf  & \valFOout & \valFOhosp\% & \valFOpib \\
Deserto pleno           & \valBOn & \valBOpop  & \valBOiso & \valBOf  & \valBOout & \valBOhosp\% & \valBOpib \\
Longe, conectado        & \valDOn & \valDOpop  & \valDOiso & \valDOf  & \valDOout & \valDOhosp\% & \valDOpib \\
Sem deserto             & \valNEn & \valNEpop  & \valNEiso & \valNEf  & \valNEout & \valNEhosp\% & \valNEpib \\
\bottomrule
\end{tabular}
\par\smallskip
\begin{minipage}{0.92\linewidth}
\footnotesize
\textit{Nota:} Medianas municipais, exceto N e população (soma). ``Fluxo que
sai'' $=$ fração das internações em hospital fora do município. ``Hosp.\ local''
$=$ \% de municípios com ao menos um leito hospitalar instalado e ativo na
janela. Cobertura privada (ANS) mediana, omitida da tabela por brevidade:
\valFOans\%, \valBOans\%, \valDOans\% e \valNEans\%, na ordem das linhas. Fonte:
SIH/DATASUS, CNES, IBGE, ANS.
\end{minipage}
\end{table}

\subsection{Quatro municípios, por inspeção}

Os quadrantes ganham concretude em casos individuais. \textbf{Fonte Boa (AM)}
está a \valFonteBoaIso~km do hospital mais próximo pela malha rodoviária, mas
\valFonteBoaLocal\% de suas internações ocorrem localmente e seu burden efetivo é
de apenas \valFonteBoaF~km: o mapa o declara deserto profundo; o fluxo, não.
\textbf{Soledade (PB)} é o espelho: o hospital mais próximo está a
\valSoledadeIso~km, mas \valSoledadeOut\% dos pacientes saem e viajam, em média
efetiva, \valSoledadeF~km --- um deserto de fluxo que a distância classifica como
atendido. \textbf{Anchieta (SC)} leva o ponto ao extremo: \emph{sedia} um
hospital (distância zero) e, ainda assim, todas as suas internações deixam o
município; ter um leito não é ter o cuidado de que se precisa.% src: casos canônicos
No contraponto, \textbf{Sobral (CE)} --- polo regional do interior --- retém
seus pacientes (apenas \valSobralOut\% saem, burden de \valSobralF~km) apesar de
não ser capital: estar no interior não é, por si, ser deserto, desde que o
município seja um hub de fato. A distinção que separa Soledade de Sobral não
aparece em quilômetros; aparece no fluxo.
